\chapter{Conclusion}
\label{ch:conclusion}

% Restate the research question: can AMR-based static analysis help detect prompt injection?
% Answer: yes, within a well-defined scope — explicit system policies, AgentDojo-style
%   agent environments — and with a formally motivated detection class.

% Summarize the key contributions:
%   - Scoped formal definition of prompt injection as AMR-graph contradiction.
%   - Static firewall (Rule1+24) with zero false positives and high detection rate.
%   - Demonstrated that structural graph similarity (smatch) is insufficient alone.
%   - Authorization annotation mechanism for the data/instruction boundary.

% Restate the honest scope: broad-delegation attacks remain outside the defined class.
%   This is a clean boundary, not a gap — it identifies the preconditions under which
%   AMR-based static detection provides formal guarantees.

% Close with the broader implication: AMR offers a principled, injection-safe
%   representation layer for policy-constrained LLM agents.
%   The abstraction mismatch extension and taint-tracking directions suggest a path
%   toward wider coverage without sacrificing the injection-safety property.
